\subsection{Network Architectures and Expressivity}

The migrated deep-feedforward-network notes describe neural networks as learned feature maps followed by prediction heads.

\begin{definition} \label{def:sl-feedforward-network} \hlt{Deep Feedforward Network}\\
A deep feedforward network represents a map
\[
f(x;\theta)=f_L\circ f_{L-1}\circ\cdots\circ f_1(x),
\]
where typical layers have the form
\[
h^{(k)}=g^{(k)}(W^{(k)}h^{(k-1)}+b^{(k)}).
\]
Training chooses parameters $\theta$ to make this map small-risk for a task.
\end{definition}

\begin{remark} \label{rem:sl-learned-feature-map} \hlt{Learned Feature Maps}\\
Classical nonlinear prediction fixes a feature map $\phi(x)$ and learns a linear head.  Deep learning learns both the representation $\phi(x;\theta)$ and the head, usually losing convexity but gaining adaptable features.
\end{remark}

\begin{definition} \label{def:sl-output-and-hidden-units} \hlt{Output and Hidden Units}\\
Linear output units model conditional means, sigmoid units parameterize Bernoulli probabilities, and softmax units parameterize distributions over finitely many classes.  Hidden units commonly use rectified linear, leaky rectified, maxout, hyperbolic tangent, or other nonlinear activations.
\end{definition}

\begin{theorem} \label{thm:sl-universal-approximation-form} \hlt{Universal Approximation Theorem Form}\\
A feedforward network with a linear output layer and at least one sufficiently wide hidden layer with a suitable non-polynomial activation can approximate every continuous function on a compact subset of $\R^d$ arbitrarily well in the uniform norm.
\end{theorem}

\begin{remark} \label{rem:sl-depth-expressivity} \hlt{Depth and Expressivity}\\
Universal approximation is an existence result, not a training guarantee.  Depth can represent some compositional and piecewise-linear functions with far fewer units than a shallow network.
\end{remark}
